\chapter{结论与展望}
\label{cha:conclusion}

\section{论文工作总结}

本文围绕“基于联邦学习的车联网入侵检测系统设计与实现”这一题目，结合当前仓库中的项目代码、实验文档、运行产物与学校论文模板，完成了一套本科毕业论文初稿的正文组织，并对系统的设计思路、关键实现和实验结果进行了系统整理。与仅停留在算法描述层面的研究不同，本文特别强调工程闭环的完整性，即从原始数据预处理、联邦训练调度、模型评估到结果报告生成，都能够在现有代码中找到对应实现与输出证据。

在系统设计方面，本文以 `main.py` 为统一入口，将系统划分为预处理、训练和测试三种命令行模式，并通过 `config.py` 中的配置模型对数据路径、训练轮数、客户端数量、节点选择权重、Top-K 比例和量化位宽等关键参数进行统一管理。与此同时，项目还引入了随机种子设置、日志输出、配置快照与运行元数据保存机制，使得实验过程具备了较好的可复现性和可追踪性。这些设计不仅满足了车联网入侵检测原型的基本运行要求，也为论文写作与后续扩展打下了较清晰的系统基础。

在系统实现方面，本文详细分析了数据预处理模块、客户端切分与状态建模机制、联邦训练与聚合流程、参数压缩与通信统计方法，以及模型评估与证据报告落盘方式。当前系统采用轻量化 CNN-LSTM 作为检测模型，在联邦训练阶段实现了基于算力、电量和信道质量的动态节点选择机制，并结合轻量公平性惩罚缓解频繁选点问题；在模型上传阶段，则通过 Top-K 稀疏化与定点量化减少通信负担。整体上，项目已经形成了“预处理—联邦训练—本地评估—报告生成”的闭环原型，能够作为车联网联邦入侵检测系统的工程化基础实现。

在实验验证方面，本文基于仓库现有的改进方案与标准 FedAvg 基线方案结果进行了对照分析。实验结果表明，在当前演示数据和小规模验证场景下，改进方案能够在 Accuracy、Precision、Recall 和 F1-score 保持不变的同时，使总训练时间下降约 12.27\%，平均轮次耗时下降约 8.51\%，总通信量下降约 80.28\%，平均隐私代理分数也高于基线方案。这说明动态节点选择、Top-K 压缩与量化机制在当前工程条件下具备较明显的综合收益，尤其在通信效率方面表现突出。

与此同时，本文也明确指出了当前系统和实验结果的边界。现有评估样本量较小，False Positive Rate 偏高，且尚未在真实规模的 VeReMi 或 Car-Hacking 数据集上完成多轮重复实验，因此现阶段结论更适合用于证明系统原型的可运行性与方法链路的可分析性，而不能直接外推为真实车联网环境中的最终性能结论。换言之，本文的主要贡献在于完成了一套事实可核验的系统设计与实现整理，并给出了后续继续深化实验与完善文献的可靠起点。

\section{工作展望}

面向后续工作，本文认为至少可以从以下几个方面继续深化。第一，应尽快引入真实的 VeReMi、Car-Hacking 或同类车联网日志数据集，对现有系统进行多次重复实验，并报告均值、标准差及必要的显著性分析，以增强实验结论的可信度。第二，可围绕 FedProx、多种客户端参与率、多组 Top-K 比例与量化位宽，构造更完整的对照实验矩阵，从而更系统地分析不同改进策略之间的相互作用。第三，可进一步改进模型结构与特征工程，例如结合更细粒度的时序特征、异常模式特征或注意力机制，以降低当前小样本场景下误报率偏高的问题。

此外，若希望系统更接近真实车联网部署环境，还可以在未来工作中补充更真实的节点状态采集机制、在线数据流处理能力和更严格的隐私保护手段。例如，可将当前的隐私代理指标扩展为差分隐私或安全聚合等具备理论保障的方法；也可以将单机顺序模拟的联邦训练原型进一步扩展为多进程甚至多节点的真实联邦实验平台。对于论文写作本身，后续仍需结合开题报告、任务书和正式参考文献，对绪论中的研究现状、文献综述和题目背景进行进一步润色与补强，从而形成更加完整、更加符合正式答辩要求的毕业论文终稿。
